% SHARD-ID: LB-15-NEGATIVE-RESULTS
% SHARD-TITLE: Negative results — refutations, obstructions, and dead ends
% SHARD-SOURCES: claims/CLAIMS.md (all REFUTED rows: A2-orbit-r1, G0-soft-r1,
%   S2, ML5, M, SPT-B-r1, SPT-nogo, M-INDEX-LA-strong); theory/amp.md;
%   theory/dh-hunt.md; theory/dh-semiclassical.md; theory/cat-hunt.md;
%   theory/soft-2d-hunt.md (NO-CAT-SOFT, §3 and the honest-status list);
%   theory/corner-a-kinks.md (1.9); theory/corner-a-goldstone.md (1.6.2.7);
%   theory/ml5-universality.md (1.4); theory/corner-b-draft.md (§10);
%   theory/spt-rebuild.md (1.3.2.3, 1.4); theory/memory-index.md (1.3, 1.12);
%   theory/verdicts/w21-joint-critic.md (W6-O2, item 6); HANDOFF.md
% SHARD-CLAIMS: the refutation ledger (A2-orbit-r1, G0-soft-r1, S2, ML5, M,
%   SPT-B-r1, SPT-nogo, M-INDEX-LA-strong — summary entries only, detailed
%   environments owned by the sections cross-referenced below) and the three
%   negatives owned here: the amputation obstruction (AMP / AMP-OBS, D24(b),
%   D24(d)), the Duistermaat--Heckman hunt and its semiclassical survivor,
%   and the no-categorical-soft-theorem fence (NO-CAT-SOFT) with its
%   operator/categorical split.

\section{Negative results: refutations, obstructions, and dead ends that shaped the campaign}
\label{sec:negative-results}

% --- shard-local semantic macros (not in main.tex's preamble) -------------
\newcommand{\Irrsimple}{\operatorname{Irr}}    % simple objects of a fusion category
\newcommand{\Funcat}{\operatorname{Fun}}       % category of module functors
\newcommand{\Dendp}{\mathcal{D}_{\mathcal{M}}} % endpoint category on register M

The campaign's constitution treats a sharp refutation, delivered together with
the weaker statement that survives it, as a result in its own right: a pile of
suggestive analogies is worth less than one clean demonstration that a
tempting statement is false. That rule has been applied often enough that the
negative outcomes now form a body of knowledge of their own, and this chapter
collects them in one place so that no reader — and no later session — walks a
dead end twice.

Two kinds of material appear here. The first is a ledger of every withdrawn
claim: the statement as it was once made, the mechanism that killed it, and
what remains true. Each of those claims belongs to a positive chapter, where
it is presented in full alongside the repaired result it forced; the entries
below are two-to-four-sentence summaries with cross-references, so that this
chapter can be read on its own as the campaign's ``what is false'' list. The
second is the set of negative results owned by this chapter, which have no
positive home: an obstruction theorem about the soft-leg amputation constant,
a fully mapped and closed search for an equivariant-localisation mechanism,
and a no-go fence separating what a fusion category can and cannot fix.

One procedural point must be stated once, because it governs the provenance
blocks throughout. Under the campaign's process rules a negative result earns
no adversarial review rounds: it is recorded, its checker (if any) is run, and
the line is closed. Consequently the ``reviewed in'' field of every statement
below is empty, and none of these statements carries a \statusProved tag. The
absence of a review artifact is deliberate and is not a defect; where a
verdict file does exist for a neighbouring positive row it is named, and
nowhere else.

\subsection{The refutation ledger}
\label{sub:refutation-ledger}

\paragraph{The vacuum-pair orbit label.} \statusRefuted{} The first form of
the asymptotic-symmetry orbit statement asserted that the set of vacuum pairs
is the orbit $\A = (G_L \times G_R)/G_{\mathrm{diag}}$, labelled by the class
$[g_L g_R^{-1}]$. This is false: $G_L \times G_R$ acts componentwise on
$\Omega_{\mathrm{vac}} \times \Omega_{\mathrm{vac}}$, so the stabiliser of a
pair $(\alpha',\beta')$ is $H_{\alpha'} \times H_{\beta'}$ and the space is
$(G/H_\alpha) \times (G/H_\alpha)$, not a quotient of $G \times G$ by the
diagonal; the proposed label is not invariant under the global diagonal
symmetry and was withdrawn. What survives is stronger and cleaner: modulo the
diagonal symmetry the complete invariant of a vacuum pair is the double coset
$\mathfrak{d}(\alpha_L,\alpha_R) = H_\alpha\, g_L^{-1} g_R\, H_\alpha \in
H_\alpha \backslash G / H_\alpha$, which for $G=SU(2)$, $H_\alpha=U(1)$ is the
relative angle $\cos\theta = \hat n_L \cdot \hat n_R$; the classification
additionally carries the transitivity hypothesis explicitly, since covariance
of the vacuum family does not imply it. See \cref{sec:corner-a}.

\provenance{A2-orbit-r1 (refuted); repaired statement A2(e)}%
{disproved in \texttt{theory/corner-a-kinks.md} step (1.9).(2.5).(3.3)}%
{\texttt{theory/checks/corner\_a\_check.py} C7}{---}

\paragraph{The universal soft factor and its Adler zero.} \statusRefuted{}
The first form of the Goldstone-tensor corollary called the kinematic factor
$e^{ik}-1$ a \emph{universal} soft factor, and inferred from it both an Adler
zero and the hard-independent linear coefficient of the two-magnon oracles.
Both inferences are withdrawn. In a matrix element the operator identity says
only
\begin{equation}
  \bra{\mathrm{out}}[H,Q_k]\ket{\mathrm{in}}
  = (e^{ik}-1)\,\bra{\mathrm{out}}J_k\ket{\mathrm{in}},
\end{equation}
so all dependence on the Hamiltonian, on the symmetry direction, on the vacuum
and on the hard legs sits inside $\bra{\mathrm{out}}J_k\ket{\mathrm{in}}$: if
that matrix element is $C_{\mathrm{hard}} + O(k)$ then hard data enter at
order $k$, and if it carries a $1/k$ infrared singularity there is no zero at
all. The factor is also convention-dependent — the kernel $e^{-ikx}$ gives
$e^{-ik}-1$, and reversing the current's orientation flips it — so its
independence of $H$ is the tautology that a discrete difference was factored
out \emph{after} an $H$-dependent current had been defined. The exact
finite-window summation-by-parts identity survives untouched, and what the
argument does rederive is the one-magnon oracle only. See
\cref{sec:corner-a}.

\provenance{G0-soft-r1 (refuted); surviving statement G0(c)--(e)}%
{withdrawn in \texttt{theory/corner-a-goldstone.md} step (1.6).(2.7)}{---}{---}

\paragraph{The superseded two-body label.} \statusRefuted{} This entry is
bookkeeping rather than mathematics, and is listed only because the claims DAG
records it with the same tag as a genuine falsehood. The old label for the
minimal two-magnon core was withdrawn so that a single mathematical statement
could not carry two statuses at once; the content it named is alive and proved
under its current label, and is generalised to every spin by the exact
regular-channel ratio. Nothing was retracted here except a name. See
\cref{sec:universality}.

\provenance{S2 (withdrawn label); live content S2-2body, S2-2body-S}{---}{---}{---}

\paragraph{Unrestricted soft universality.} \statusRefuted{} The claim that
the soft factorisation holds for every local or quasi-local source is false,
and the counterexample is explicit and four sites wide. With
\begin{equation}
  D := S^-_0 S^-_1 - S^-_1 S^-_2 + S^-_2 S^-_3 - S^-_0 S^-_3,
  \qquad
  O_\eta := S^-_0 + \eta D,
\end{equation}
every $O_\eta$ is local; because $D\ket{\Omega}$ lies in the two-magnon sector
it leaves the one-magnon amplitude $M_1^{O_\eta}(h)$ completely independent of
$\eta$, while the two-magnon amplitude acquires the first jet
\begin{equation}
  M_2^{O_\eta}(k,h) = \eta\bigl\{ 2i\,(1-e^{-3ih})\,k + O_I(k^2) \bigr\}.
\end{equation}
On any hard packet supported where $1-e^{-3ih} \neq 0$ the parameter $\eta$
therefore moves the linear soft coefficient without moving the hard process at
all, so no process-independent coefficient can exist on the unrestricted local
class. The surviving statements are the repaired criterion — process
independence holds exactly when the source's amputated commutator has
vanishing first jet — and the conditional factorisation on the restricted
class $\classSW(\rho)$, whose conclusion fixes a \emph{shape} and not a
number. See \cref{sec:universality}.

\provenance{ML5 (refuted); surviving statements ML5-A, ML5-B}%
{\texttt{theory/ml5-universality.md} step (1.4)}%
{\texttt{theory/checks/ml4\_check.py}, the $O_\eta$ obstruction}{---}

\paragraph{Memory as the direct-current limit of the soft factor.}
\statusRefuted{} The brief's memory conjecture, read literally — the wall
displacement equals the zero-frequency limit of the soft factor, summed over
the event — is false in the easy-axis chain. Two independent reasons kill it:
conditional charge bookkeeping fixes the same displacement coefficient for
every momentum, every anisotropy and every packet once the scattering channels
exist, whereas any soft-factor expression varies with all three; and the
displacement is entirely insensitive to the transmission phase, since a purely
transmitting wall with vanishing phase still displaces by exactly $-1/s$. The
soft factor is a phase, the memory quantum is a charge. What survives is the
exact flux statement — the displacement is the finite-time direct-current
weight of the \emph{physical} boundary current — together with the conditional
quantisation statement, into which soft data enter only through the
transmission probability $T(k)$; no reading of virtual or bond data is
claimed. See \cref{sec:memory-quantization}.

\provenance{M (refuted); surviving statements M-flux, M-quant (conditional)}%
{\texttt{theory/memory-quantization.md} §1; \texttt{theory/corner-b-draft.md} §§2, 6, 10}%
{\texttt{theory/checks/mquant\_check.py} (flux residues plus an empirical scan only)}{---}

\paragraph{Pointwise blindness of closed adjoint contractions.}
\statusRefuted{} An early symmetry-protected-phase argument claimed that
closed contractions built only from the adjoint action $\Ad(V)$ of the virtual
symmetry are pointwise blind to the cohomology class $[\omega]$. They are not.
For the two-element group with Pauli-projective virtual action the adjoint
representation decomposes into four distinct characters, with multiplicities
$(1,1,1,1)$, whereas the scalar-trivial lift gives four copies of the trivial
character, $(4,0,0,0)$; the closed scalar $\operatorname{Tr}\Ad(V(R_x))$ is
then $0$ in the first case and $4$ in the second. What survives is the exact
multiplier cancellation for ordered closed insertions, where the endpoint
action is $V(g) \otimes \bar V(g)$ and the projective phases cancel identically,
together with the continuity statement for normalised coefficients — which
becomes topological only after a separate local-constancy proof, since
continuity does not imply local constancy. See \cref{sec:spt}.

\provenance{SPT-B-r1 (refuted); surviving statements SPT-B-mult, SPT-B'}%
{disproved in \texttt{theory/spt-rebuild.md} step (1.3).(2.3)}%
{\texttt{theory/checks/spt\_rebuild\_check.py} S-C2}{---}

\paragraph{The all-orders cohomology no-go.} \statusRefuted{} The companion
early claim was that the cohomology class can appear in no coefficient at all,
an edge residue included. This over-read Whitehead's lemma, which trivialises
only the \emph{infinitesimal} central cocycle of a compact semisimple Lie
algebra and only in a stated phase section; it never erases the dimension, the
weights, the Casimir data or the integrability class of a representation, so
global edge weights and module dimensions survive. What remains true is the
narrow version: in the registered endpoint the infinitesimal cocycle is a
coboundary and is gauged away in the stated phase convention, while the edge
residue itself is the centred, phase-gauge-invariant charge whose spectrum
lies in a fixed coset of $\Z$ and whose irreducible module has dimension at
least $d_\omega$. See \cref{sec:spt}.

\provenance{SPT-nogo (refuted); surviving statements SPT-D', SPT-E'}%
{\texttt{theory/spt-rebuild.md} step (1.4)}%
{\texttt{theory/checks/spt\_rebuild\_check.py} S-C4}{---}

\paragraph{A sector-wide regularised total charge.} \statusRefuted{} It was
claimed that the kink hypotheses together with integrality force the existence
of a self-adjoint regularised total charge operator in every kink
representation. They do not. The decisive counterexample is a product state in
a kink sector whose on-site deviations decay only polynomially: with
$\varepsilon_n = (n+1)^{-1/2}$ the accumulated variance
$V_N = \sum_{n \le N}\varepsilon_n^2$ diverges, the window-charge distribution
spreads with $\sup_k P(L_N = k) \le C V_N^{-1/2}$, and hence
$\braket{\Omega_\varrho, (\Qhat_{W_N,c_0} - i)^{-1}\Omega_\varrho} \to 0$,
which no self-adjoint limit can produce, since for any self-adjoint $\Qhat$
the imaginary part of that resolvent expectation is strictly positive. The
mean-controlling hypothesis governs first moments, not charge fluctuations,
and hence not implementability. A second, on-folium mechanism reaches the same
conclusion for a nonscalar virtual circle action, but it is a strengthening
under strictly more hypotheses, not an independent proof of the same
statement. What survives is that the relative window charge exists as a
\emph{limit vector} on states with exponentially clustering tails, where it
coincides with the exactly conserved first-moment wall coordinate; no
sector-wide operator is constructed anywhere in the corpus. See
\cref{sec:memory-index}.

\provenance{M-INDEX-LA-strong (refuted); surviving statements LD-ID, ACE-LD-eps}%
{\texttt{theory/memory-index.md} step (1.3) (fluctuation counterexample, off-folium; alone sufficient) and step (1.12) (nonscalar virtual circle, on-folium; a strengthening under additional hypotheses)}%
{\texttt{theory/checks/memory\_index\_check.py} IDX-C3, IDX-C8}{---}

\begin{remark}[Partial refutations that live inside proved rows]
Three further negative findings are not listed above because they are
corrections \emph{within} statements that are otherwise proved, and are
presented with those statements rather than here: the two-magnon Ward
projection erratum, where a scalar-normalised projection formula holds at one
magnon only and its register-dependent correction is non-scalar at two or more
(\cref{sec:wave-operators}); the withdrawal of the raw half-line displacement
formula in favour of the leg-subtracted event identity $2s\,\delta x +
(q_{\mathrm{out}} - q_{\mathrm{in}}) = 0$ (\cref{sec:memory-quantization});
and the refutation of the fully polarised product-state form of the tail vacua
under the general kink hypotheses, which survives only when the tail density is
a simple on-site eigenvalue (\cref{sec:local-relaxation}). Each is recorded as
an erratum next to the theorem it scopes.
\end{remark}

\subsection{The amputation obstruction: the class fixes only a square root}
\label{sub:amputation-obstruction}

The remainder of this chapter presents the negatives that have no positive
home. The first is the sharpest of them, because it names precisely what a
future proof must supply.

\begin{definition}[Soft-leg amputation constant]
\label{def:amputation-constant}
Work at a fixed tail density $\rho > 0$, put $Z_\rho := 2\rho$, and assume the
restricted source class $\classSW(\rho)$ is nonempty and the asymptotic
one-magnon kernel exists. Fix the normalisation convention in which the hard
leg is amputated identically in the one-magnon amplitude $M_1^O$ and the
two-magnon amplitude $M_2^O$, while the extra soft leg of $M_2^O$ is a
unit-weight, delta-normalised asymptotic magnon. Decompose $M_2^O$ into a
descendant term $E^O_{\mathrm{desc}}$, an orthogonal-current term and a
direct/contact term, the last two bounded by $O(k^2)$; fix the descendant
current residue in the charge-created normalisation to $2iv_h M_1^O(h)$; and
for $k \neq 0$ define the descendant quotient
\begin{equation}
  L(k,h) := \frac{E^O_{\mathrm{desc}}(k,h)}{(e^{ik}-1)\,2iv_h\,M_1^O(h)},
\end{equation}
assumed to have a process-independent, uniformly $C^1$ extension to $k=0$.
The \emph{soft-leg amputation constant} $\aleg(\rho) \neq 0$ is then defined by
the value of that extension,
\begin{equation}
  L(0,h) = \aleg(\rho)\,\frac{-i\chi(h,0)}{v_h},
  \label{eq:aleg-defn}
\end{equation}
with $\chi$ the channel sign. In this convention the conditional
factorisation theorem reads $M_2^O(k,h) = 2i\,\aleg(\rho)\,\chi\,k\,M_1^O(h)
+ O_{L^2(I)}(k^2)$, which contains no number at any density.
\end{definition}

\provenance{D24(b), D24(d)1--3, D24(d)3b; consumers ML5-B, D24-VAL, AMP}%
{\texttt{definitions.md} D24; \texttt{theory/amp.md} §1}%
{\texttt{theory/checks/d24d3\_normalization\_check.py} D24N-C1--C3 (kinematics of the conclusion, not its constant)}%
{\texttt{theory/verdicts/d24d3-adjudication-r5.md} §5 is the merged text of record for the clause repair}

The natural hope was that this constant is fixed by the definition of the
universality class, with the physically expected value $\aleg(\rho) = 1/Z_\rho$
following from a careful accounting of how a charge-created soft leg is
converted into an asymptotic one. That hope is refuted at the level of
derivability.

\begin{refutedclaim}[The universality class fixes the amputation constant]
\label{ref:amputation}
Withdrawn statement: the normalisation conventions and analytic relations
defining the class $\classSW(\rho)$ determine the value of $\aleg(\rho)$, and
in particular force $\aleg(\rho) = 1/Z_\rho$. This is false. Those relations
determine $\aleg(\rho)$ only up to the constraint $\aleg(\rho) \neq 0$
together with process independence; on the charge-created reading they supply
exactly one factor $Z_\rho^{-1/2}$, and the remaining scalar is free. The
statement is neither proved nor physically refuted — it is undecidable from
the present hypotheses, and a completed proof requires a genuinely new
dynamical input, displayed in \eqref{eq:amp-missing} below.
\end{refutedclaim}

\emph{Mechanism.} Three independent steps close the question. The first is an
exact leg-conversion computation. On a fully polarised tail with $Z_\rho = 2S$,
the charge-created soft leg on the vacuum and the descendant vector
$u_N := Q^-_q\ket{h}$ built on a hard magnon at a commensurate momentum satisfy,
on a finite ring,
\begin{equation}
  Q^-_k\ket{\Omega} = \sqrt{Z_\rho}\,\ket{k},
  \qquad
  \norm{u_N}^2 = Z_\rho N - 2 .
  \label{eq:leg-conversion}
\end{equation}
The rank-one orthogonal projection onto the descendant line must then be read
in the normalised register:
\begin{equation}
  P_N w_N
  = u_N \frac{\braket{u_N,w_N}}{Z_\rho N - 2}
  = \frac{u_N}{\sqrt{Z_\rho N - 2}}\;
    \frac{\braket{u_N,w_N}}{\sqrt{Z_\rho N - 2}} .
  \label{eq:amp-projection}
\end{equation}
The first square root normalises the \emph{vector}; only the second is the
amplitude coefficient against that normalised vector. Counting both as
amplitude factors changes register mid-equation, and that mutation is
registered as an explicit failing mode of the checker. Removing the common
finite-ring delta-normalisation leaves the exact descendant conversion factor
$(Z_\rho - 2/N)^{-1/2}$, so that
\begin{equation}
  \lim_{N\to\infty}\bigl(Z_\rho - 2/N\bigr)^{-1/2} = Z_\rho^{-1/2},
  \qquad
  \bigl(Z_\rho - 2/N\bigr)^{-1/2}
  = Z_\rho^{-1/2}\Bigl[1 + \frac{1}{Z_\rho N} + O(N^{-2})\Bigr].
\end{equation}
The finite-size mismatch is therefore a vanishing correction, not a surviving
second factor.

The second step checks the four operations that might have supplied a second
factor, and finds that none does. The descendant projection is already
exhausted by \eqref{eq:amp-projection}. Windowed-charge smearing is linear and
supplies no density-dependent scalar, since the packet register fixes
delta-normalised kernels and soft envelopes with value $1$ at $k=0$; a
constant envelope $Z_\rho^{-1/2}$ would change that value and would be a
forbidden second leg rescaling rather than a new mechanism. The two
normalisation anchors remove a rescaling freedom and impose the Ward residue
as a membership condition, but neither evaluates $L(0,h)$. And the finite-ring
mismatch vanishes, as just displayed. The one remaining place where a new
mechanism could live is the descendant propagation itself, which the class
definition asserts to exist and to have a $C^1$ quotient without ever
specifying a map whose normalised residue could be evaluated.

The third step makes the underdetermination explicit at the level of the
relations. Fix a compact hard window on which $v_h \neq 0$ and the channel
sign is constant, choose a nonzero profile $m \in L^2(I)$, and for any
$a \in \C\setminus\{0\}$ consider the source-linear package
\begin{align}
  M_1(h) &= c\,m(h), &
  R(h) &= c\,2iv_h\,m(h), &
  L_a(k,h) &= a\,\frac{-i\chi}{v_h}, \notag\\
  E_{\mathrm{desc},a}(k,h) &= (e^{ik}-1)\,L_a(k,h)\,R(h), &
  M_{2,a} &= E_{\mathrm{desc},a}, &
  &\text{other terms } = 0 .
\end{align}
Every clause of the class definition holds exactly: the decomposition is
trivial, the remainder bounds hold with zero remainder, the Ward residue is
$R = 2iv_h M_1$, and the descendant quotient is exactly $L_a$, with
$iv_h L_a(0,h)/\chi = a$. Two distinct values of $a$ leave $M_1$, the Ward
residue, the packet measure and the leg convention unchanged while changing
$L$ and $M_2$, so they are not related by any normalisation covariance. Hence
every nonzero constant is admitted by the displayed relations.

\emph{What a completed amputation proof must supply.} Granting the
charge-created reading of the descendant term, factor the class datum as
\begin{equation}
  \aleg(\rho) = Z_\rho^{-1/2}\,\bigl[\sqrt{Z_\rho}\,\aleg(\rho)\bigr] ,
\end{equation}
where the leg conversion \eqref{eq:leg-conversion} fixes the first factor and
the class definition constrains the bracket only to be nonzero and
process-independent. The desired value is therefore equivalent to the single
additional equation
\begin{equation}
  \sqrt{Z_\rho}\,\aleg(\rho) = Z_\rho^{-1/2} .
  \label{eq:amp-missing}
\end{equation}
Any proof of \eqref{eq:amp-missing} must be a post-conversion
descendant-propagation or LSZ residue theorem, stated after both sides are
already expressed against the same delta-normalised soft leg; it is not a leg
normalisation, and no clause of the present definition supplies it. Every
candidate argument must be checked against the exact leg-conversion
computation \eqref{eq:leg-conversion}, or it double-counts the leg.

\begin{scope}[What is and is not decided]
\label{scope:amputation}
The obstruction is at relation level: it exhibits packages satisfying the
displayed algebraic and analytic relations for every nonzero constant, and it
does not construct a microscopic local source, so it is not a physical
counterexample and does not refute the value $1/Z_\rho$ in a future nonempty
class. A separate, strictly conditional falsifier does bear on the
leg-normalisation-only value: assuming both the charge-created reading of the
descendant term and the jet-identification bridge that equates the class
multiplier's jet with the two-body physical phase jet, a pure-leg proof
predicts the jet $2/\sqrt{2S}$, whose scaled deviations against the
ansatz-free two-magnon data are $0.0004$, $0.4158$, $0.7347$ and $1.0033$ at
$S = 1/2, 1, 3/2, 2$, i.e.\ margins of $5.2$, $9.2$ and $12.5$ times the
pre-registered band of $0.08$ at the three non-degenerate spins. Neither
antecedent is proved, and without them a two-magnon phase slope and an
amputation constant are not comparable at all; the relation-level obstruction
above uses neither. Finally, both candidate value rows remain vacuous or
unknown in a further sense: nonemptiness of the class is open at every
density, so they constrain future work — any member ever exhibited must carry
the stated value — rather than supplying evidence now.
\end{scope}

\provenance{AMP \statusConjecture; the obstruction is recorded as the proposed row AMP-OBS; inputs D24(b), D24(d)}%
{\texttt{theory/amp.md} steps (1.1)--(1.4); the honest disposition is step (1.4).(2.2)}%
{\texttt{theory/checks/amp\_check.py} AMP-C1--C3 with the registered failing modes \texttt{-{}-red-double-count} and \texttt{-{}-red-pure-leg}; AMP-C4 is the conditional falsifier only; \texttt{theory/checks/d24d3\_normalization\_check.py} D24N-C8 (leg-conversion constant), D24N-C6 \texttt{-{}-red-halfpower}}%
{--- (no review round: negative results earn none by design; the claims DAG keeps AMP at conjecture)}

\subsection{The Duistermaat--Heckman hunt, fully mapped}
\label{sub:dh-hunt}

The second negative is a closed search. Its target was attractive enough to be
worth stating precisely: the memory law produced by the two-time measurement
protocol is an atomic probability measure on $\Z$, and exact integrality
statements of that kind are exactly what equivariant localisation delivers in
symplectic geometry. If the Duistermaat--Heckman theorem — exactness of
stationary phase for a Hamiltonian torus action, equivalently the statement
that the pushforward of Liouville measure under the moment map is piecewise
polynomial — could be applied to the outcome law, the campaign would gain both
the support and the weights from geometry alone. It cannot, and this
subsection records why in enough detail that the route need not be reopened.

\begin{definition}[The two objects that were to be identified]
\label{def:dh-objects}
The \emph{Duistermaat--Heckman law} of a compact Hamiltonian $T$-manifold
$(M,\varpi,\mu)$ is the pushforward $\mu_*(\mathrm{vol}_\varpi)$ of Liouville
measure by the moment map: it depends on the manifold, its symplectic form,
the torus action and the moment map, and on nothing else. The \emph{memory
outcome law} is $p_W = \sum_{\nu \in \Z} p_W(\nu)\,\delta_\nu$, where each
weight is a matrix element of a product of two spectral projections of the
same windowed charge at two different times, followed by an ordered window and
Ces\`aro limit; it depends on the state, on the dynamics, on the back-action
of the first measurement, and on the limit order.
\end{definition}

\begin{refutedclaim}[The memory outcome law is a Duistermaat--Heckman measure]
\label{ref:dh-law}
Withdrawn statement: the escaped-charge law of the memory protocol is the
Duistermaat--Heckman measure of a natural Hamiltonian torus space, so that its
integer support and its weights follow from equivariant localisation. No
candidate identification survives. The two objects have different types; the
one construction that produces atoms does so only by inserting the desired
atoms and their volumes by hand; and every route that reaches the soft
coefficient consumes the representation-theoretic or scattering data it was
meant to produce.
\end{refutedclaim}

\emph{Mechanism, candidate by candidate.} Six natural routes were worked to
their stopping points.

\begin{center}
\begin{tabular}{@{}p{0.30\textwidth}p{0.62\textwidth}@{}}
\toprule
route & why it stops \\
\midrule
the memory law as a localisation measure &
a Duistermaat--Heckman measure on a positive-dimensional Hamiltonian torus
manifold is a continuous, piecewise-polynomial pushforward, while the outcome
law is atomic and state- and dynamics-dependent \\[2pt]
weight-multiplicity asymptotics &
the projective law takes the weight multiplicities as its Dirichlet
parameters, so it consumes them rather than producing them \\[2pt]
localisation for the soft coefficient &
the internal circle acts on a one-magnon band by a scalar character, so the
whole band is fixed and $k_s = 0$ is not an isolated fixed locus \\[2pt]
wall crossing as phase rigidity &
localisation walls are critical moment values inside one fixed space, whereas
the phase-diagram sweep changes the Hamiltonian, the tensor and the symplectic
data themselves \\[2pt]
convexity and integral vertices &
integral vertices do not make interior moment values integral: already on the
one-qubit sphere the moment map fills $[-\tfrac12,\tfrac12]$ \\[2pt]
the kink K\"ahler cylinder &
in two dimensions the moment-map identity alone gives the answer, so
localisation performs no inference \\
\bottomrule
\end{tabular}
\end{center}

The second route deserves its exact computation, because it kills the premise
cleanly. Take two spin-$\tfrac12$ sites, $V = (\C^2)^{\otimes 2}$, with circle
generator $Q = S^z_1 + S^z_2$, so that $V = V_{-1} \oplus V_0 \oplus V_{+1}$
with dimensions $(1,2,1)$; let $M = \mathbb{P}(V) = \mathbb{CP}^3$ carry
normalised Fubini--Study measure and the moment map
$\mu([v]) = \braket{v,Qv}/\braket{v,v}$. The normalised weight-counting law and
the localisation law are then
\begin{equation}
  \eta_{\mathrm{wt}}
  = \tfrac14\delta_{-1} + \tfrac12\delta_{0} + \tfrac14\delta_{+1},
  \qquad
  \mathrm{d}\eta_{\mathrm{DH}}(x)
  = \tfrac32\,(1-\abs{x})^2\,\mathbf{1}_{[-1,1]}(x)\,\mathrm{d}x ,
  \label{eq:dh-cp3}
\end{equation}
the second obtained by writing the block fractions of a Fubini--Study random
ray as a Dirichlet vector with parameters $(1,2,1)$ and integrating out the
neutral block. They are not the same measure, and not merely by normalisation:
the first is supported on three points and the second is absolutely
continuous, and their second moments are $\tfrac12$ and $\tfrac{1}{10}$
respectively. The general statement is the same in every torus
representation $V = \bigoplus_\lambda V_\lambda$: the block fractions
$(p_\lambda)$ are Dirichlet with parameters $\dim V_\lambda$ and
$\mu_T = \sum_\lambda p_\lambda \lambda$, so the density is a spline whose
inputs already include every multiplicity.

The third route fails for a structural reason worth stating on its own. The
internal circle acts on a one-magnon state by its charge character, which is
projectively trivial, so the entire one-magnon band is fixed rather than a
point; on a finite ring the translation symmetry is the cyclic group and has
no infinitesimal generator; and in infinite volume the Brillouin circle labels
irreducible translation characters and is not a Hamiltonian circle action on a
compact phase space. The two meanings of ``stationary phase'' also fail to
meet: localisation needs a phase that is a moment map, while the scattering
estimates use $t\omega(k) - kx$ plus an on-shell scattering phase, and no
definition in the corpus turns a dispersion relation into an internal-charge
moment map. Finally, the soft coefficient is not a symmetry-only integral at
all: multiplying the Hamiltonian by a constant leaves the projective symmetry
geometry unchanged while rescaling the current-velocity residue $2iv_S(h)$.

The sixth route is instructive because it \emph{succeeds} and thereby shows
what success would be worth. On a two-dimensional Hamiltonian circle manifold
with coordinates $(x_0,\varphi)$, invariant form $\varpi = \Omega(x_0)\,
\mathrm{d}x_0 \wedge \mathrm{d}\varphi$ and moment map $\mu(x_0)$, the defining
identity $\mathrm{d}\mu = \iota_{\partial_\varphi}\varpi$ gives
$\mu'(x_0) = -\Omega(x_0)$, hence a constant pushforward density and
\begin{equation}
  \int_{x_0}^{x_0+1}\!\!\int_0^{2\pi} \varpi
  = 2\pi\bigl[\mu(x_0) - \mu(x_0+1)\bigr] .
\end{equation}
With the calibration $\abs{\Delta\mu} = 2s$ the area per lattice translation is
$4\pi s$ and the position quantum is $1/(2s)$ — the right answer, obtained by
integrating the defining identity and nothing more. That is geometric
interpretation, not inference, and it was pre-registered as the
reproduction-only failure mode before the hunt began.

\begin{observation}[The coadjoint-orbit soft-phase envelope]
\label{obs:sc-envelope}
One genuinely new result did survive the hunt, on the phase rather than on the
weights. Write the classical spin phase space as a product of coadjoint orbits
$P_S = \prod_x \mathcal{O}_S$ with $\mathcal{O}_S \cong S^2$ and form
$\Omega_S = S\sum_x \omega_{\mathrm{KKS},x}$, so that the effective
semiclassical parameter is $1/S$. Put $p = (k_s+k_h)/2$, $q = (k_s-k_h)/2$,
$c_p = \cos p$, $c_q = \cos q$, $s_q = \sin q$, and in the regular chamber
$c_p s_q \neq 0$ set
\begin{equation}
  F = \frac{c_q\,(c_p - c_q)}{c_p\,s_q},
  \qquad
  d = -1 + \frac{c_q}{c_p} .
\end{equation}
Then the exact two-magnon contact phase, taken on the continuous branch with
$\delta_S(0,k_h) = 0$, is
\begin{equation}
  \delta_S(k_s,k_h) = 2\arctan\!\Bigl(\frac{F}{2S+d}\Bigr)
  = \frac{F}{S} + \frac{G}{S^2} + \frac{H}{S^3} + O(S^{-4}),
  \qquad
  G = -\frac{Fd}{2},
  \quad
  H = \frac{Fd^2}{4} - \frac{F^3}{12},
  \label{eq:sc-envelope}
\end{equation}
with every function fixed before any comparison and no fitted coefficient
anywhere. At $k_s = 0$ one has $c_p = c_q$, hence $F(0,k_h) = 0$,
$\partial_{k_s}F(0,k_h) = 1$ and $d(0,k_h) = 0$, so $G = O(k_s^2)$ and
\begin{equation}
  \partial_{k_s}\delta_S(0,k_h) = \frac{1}{S} :
\end{equation}
the proved finite-spin soft slope is exactly the leading semiclassical phase,
and the first subleading term cannot renormalise the linear soft coefficient.
\end{observation}

\emph{Status and evidence.} This is a speculative result about the large-spin
envelope of the two-magnon phase, not a second proof of the finite-spin slope:
\eqref{eq:sc-envelope} reparameterises the exact contact ratio that the proved
two-body theorem already supplies, and its new content is that the hierarchy
is visible and testable. Eight frozen soft-window rows at $S \in \{1/2,1\}$
and $k_s \le 0.12$, previously compared only with the full exact answer, agree
with the truncation $F/S + G/S^2$ to within $2.62\%$ in the worst row — the
softest-momentum row at $S = 1/2$ — inside the declared $3\%$ descriptive
gate, and adding the parameter-free $G/S^2$ term reduces the error of every
one of them by at least $89.8\%$. An independent exact-diagonalisation
sequence on an eighteen-site ring, running through $S = 8$ and computed from
finite-volume level displacements rather than from the closed-form contact
ratio, lies inside the derived pointwise remainder bound at every spin, with
second-order residuals decreasing monotonically and $\max_S S^3\abs{\text{error}}
= 0.0089$. Identifying $G$ with any particular fluctuation determinant is
speculative and is not needed; the checked statement is only that $G/S^2$ is
the first subleading contact correction of exactly that size.

\begin{scope}[No localisation law for the outcome probabilities]
\label{scope:dh-weights}
The envelope result does not extend to the memory weights, and the negative
verdict there is sharp. The one-magnon protocol holds the escaped charge at
order one rather than letting a highest weight grow, so the usual
semiclassical bridge has the wrong limit: dividing the escaped integer by the
spin or by the window size collapses the law to a point mass at zero, while
declining to divide leaves an atomic law rather than a continuous density. The
frozen memory data confirm the split. Thirteen memory rows at
$S \in \{1/2,1,3/2\}$ have displacement-per-unit-time within $2.69\%$ of
$-1/S$, well inside the file's declared band — but that shared $1/S$ is the
kinematic conversion from one unit of transmitted charge to wall displacement,
not a geometric volume; the transmission weights themselves vary with the
anisotropy, the packet momentum, the packet quality and the trapped mass, and
neither the envelope nor any localisation volume predicts them. The
independent probe puts at least $0.99601$ of its final weight on a single
outcome with zero off-lattice mass, which tests exact support and late-time
channel concentration, again not a density. There is a real asymptotic
statement in the neighbourhood — the normalised weight multiplicities of the
line-bundle family $H^0(M,\mathcal{O}(n))$ do converge to the density in
\eqref{eq:dh-cp3} — but that is a normalised-trace multiplicity family, and
substituting it for the measured law would change the experiment; inserting
the transmission probability as an equivariant weight would insert the answer.
The absence of one further artifact should be noted explicitly: the hunt shard
deliberately has \emph{no} checker, since verification machinery is not built
for a terminal negative side-result, whereas the surviving envelope does carry
one.
\end{scope}

\provenance{no claims-DAG row (negative lane record, plus a speculative companion-paper result)}%
{\texttt{theory/dh-hunt.md} §§1--5 (the six verdicts and the exact projective computation); \texttt{theory/dh-semiclassical.md} §§1--4 (the envelope, the error model and the data comparisons)}%
{\texttt{theory/checks/dh\_semiclassical\_check.py} DHSC-C0--C4 with failing modes \texttt{-{}-red wrong-leading} and \texttt{-{}-red drop-fluctuation}; no checker exists for \texttt{theory/dh-hunt.md} by design}%
{---}

\subsection{The no-categorical-soft-theorem fence}
\label{sub:no-cat-soft}

The third negative fixes the boundary between what a fusion category can
determine and what it cannot. It matters because the two-dimensional lane's
positive selection theorem is genuinely categorical, and it would be easy to
over-read that success as a categorical soft theorem.

\begin{definition}[Endpoint register and endpoint form factor]
\label{def:endpoint-register}
Let $\mathcal{C}$ be a finite unitary fusion category represented by exact
matrix-product operators, and let $\mathcal{M}$ be the chosen indecomposable
module category recording the boundary or phase register. The category of
topological endpoint charges on that register is the dual category
$\Dendp := \Funcat_{\mathcal{C}}(\mathcal{M},\mathcal{M})$; a statement
assigning endpoint labels from $\mathcal{C}$ alone is not invariantly typed. Let
$P_a$ project onto the simple sector $a \in \Irrsimple(\Dendp)$, and call an
endpoint operator \emph{pure of type $x$} when its intertwiner is resolved into
a single simple $x$. Writing $T_x(k;w)$ for a windowed, momentum-resolved
endpoint family, its allowed block in a one-dimensional fusion space
factorises as
\begin{equation}
  F^{x}_{ba}(k;w) = f^{x}_{ba}(k;w)\,I^{x}_{ba},
  \label{eq:endpoint-formfactor}
\end{equation}
where the intertwiner $I^{x}_{ba}$ is categorical while the scalar form factor
$f^{x}_{ba}$ depends on the endpoint tensor, the window, the Hamiltonian, the
external states and the normalisation.
\end{definition}

The selection half is a theorem: a pure endpoint of type $x$ satisfies
$P_b T_x P_a = 0$ unless the fusion multiplicity $N_{xa}^{\,b}$ is positive,
and a nonzero block carries a channel vector $\mu \in \Hom(b, x \otimes a)$, so
the fine categorical datum is the resolved fusion event $(x,a\to b;\mu)$ rather
than any difference of labels; when $x$ is invertible the target is unique and
the index lives in $\Inv(\Dendp)$, while for general $x$ the maximal
group-valued shadow is the universal grading, $\abs{b}\abs{a}^{-1} = \abs{x}$.
The Ising algebra shows why the coarse statement cannot be improved: $\sigma$
has quantum dimension $\sqrt2$ and $\sigma \otimes \sigma = 1 \oplus \psi$ has
two fine targets, so there is no canonical subtraction, yet $\sigma$ still
carries a definite odd universal degree. None of this fixes a number inside an
allowed channel.

\begin{refutedclaim}[A fusion category fixes a soft coefficient]
\label{ref:no-cat-soft}
Withdrawn statement: the categorical data of an endpoint — its type, its
source and target sectors, its fusion channel, its support geometry and the
exact homogeneous pulling-through and intertwiner equations — determine the
soft behaviour of its form factor, in particular an Adler zero or a universal
linear coefficient. This is false, and sharply so: fix all of that data
together with the value $T_x(0) \neq 0$ of a nonzero analytic family near
$k=0$; then no positive-order Taylor coefficient of the allowed form factor
follows. Categorical data give selection rules, not derivatives.
\end{refutedclaim}

\emph{Mechanism.} Sector projection and every pulling-through or intertwiner
equation are linear and homogeneous in the endpoint tensor, so for any analytic
scalar $c$ with $c(0) = 1$ the family
\begin{equation}
  \widetilde T_x(k) = c(k)\,T_x(k)
\end{equation}
has the same endpoint type, the same sectors, the same channel, the same
support, the same zero-momentum operator and the same homogeneous equations.
Choosing $c(k) = e^{i\alpha k^n}$ preserves the norm and every derivative below
order $n$ while changing the $n$-th derivative by $i\alpha\,n!\,T_x(0)$, and
$\alpha$ is arbitrary; imposing norm preservation therefore still leaves an
arbitrary phase jet. Without fixing $T_x(0)$ the conclusion is even simpler:
$I^{x}_{ba}$ and $k\,I^{x}_{ba}$ obey the same homogeneous categorical
equations in an allowed one-dimensional channel. Fusion data force an exact
zero only in the \emph{forbidden} blocks of the selection rule; in an allowed
block they force neither a zero nor a nonzero constant. An exact
stabiliser-code computation makes the freedom concrete rather than formal: on a
single plaquette, dressing a string operator with a slowly varying scalar phase
profile of wavelength $1/k$ leaves the endpoint syndrome, the fusion channel
and the amplitude magnitude untouched while giving the amplitude the arbitrary
first jet
\begin{equation}
  \partial_k A_{\mathrm{ph}}(0) = i\lambda\,(s_0 + s_1 + 2c) .
\end{equation}

\emph{Why gapless input is the missing ingredient.} The failure is not an
accident of bookkeeping; it identifies precisely what ordinary soft theorems
use and this setting lacks. A massive anyon stays separated from the ground
register by a positive energy as $k \to 0$, so the limit here is a
long-wavelength limit of the chosen source rather than an energetic
soft-particle limit. What gives a soft theorem its force is a gapless on-shell
state together with a Ward identity or factorisation statement, and neither a
finite tube algebra nor a non-invertible pulling-through relation contains a
current zero mode. A fusion category with a gap can therefore be expected to
supply selection rules and no coefficient; supplying a coefficient requires
naming a gapless edge Hamiltonian and its current algebra as new input, which
is a separate work order and must not be advertised as determined by the bulk
category.

\begin{scope}[Strength of the no-go]
\label{scope:no-cat-soft}
This is a data-insufficiency theorem, not a theorem about every possible
normalisation convention: a convention can remove the scalar freedom, but the
convention is then extra kinematic data and still fixes no measurable magnitude
response. The witness varies the momentum-dependent probe family and is not a
pair of fully specified Hamiltonian or tensor-network models with one fixed
normalised protocol. No claim is made that every gapped response starts at
order $k$: a constant term, an arbitrary analytic correction, or a
convention-dependent phase is compatible with the categorical data, and the
conclusion is the absence of a category-universal coefficient, not a universal
order of vanishing. Under the negative-result rule this fence is deliberately
kept out of the claims DAG; it is carried instead inside the ``not claimed''
list of the conditional selection theorem it fences, whose statement asserts no
Taylor coefficient inside an allowed channel.
\end{scope}

\begin{observation}[The fusion datum is categorical, the Ward operator is Lie]
\label{obs:cat-lie-split}
The companion finding of the same hunt is a clean split of the campaign's
machinery by what each half actually consumes. On the categorical side stand
the sector or tube projector and its selection rule, the universal grading
shadow, and — for an \emph{invertible} on-site string, where locality permits
the conjugation $U_R^\dagger H U_R - H$ to be written as a sum of two endpoint
defects — the finite string endpoint identity. On the Lie side stand the Ward
relation $D^\dagger J\psi = R\psi$ and the invertibility and scalar ladder
value of the Gram operator $D^\dagger D$, neither of which follows from fusion
data. The witness is elementary and exact: in a two-dimensional register the
operators $D$, $J_1 = D$ and $J_3 = 3D$ all carry the same odd categorical
covariance, yet their projected vectors differ by $1:3$, and replacing $D$ by
$2D$ changes the Gram value by $1:4$ without changing any category datum. A
non-invertible string has no conjugation at all, so its pulling-through
equation continues the string-endpoint half and not the current-zero-mode half.
Accordingly the sector-label theorem is the categorical half of the campaign,
and the normalised Gram-inverse current theorem is genuinely Lie-algebraic
under the hypotheses presently known.
\end{observation}

\emph{Deliberately unpromoted.} Neither this observation nor the categorical
selection theorem behind it is proposed for the Letter. The hunt's own honest
assessment is that any corpus row derived from it would have to enter at sketch
status, because the artifact has not undergone the campaign's capped hostile
review — and by the negative-results rule it does not get one. The material is
therefore parked as companion-paper content, alongside the semiclassical
envelope of \cref{sub:dh-hunt}, pending an explicit decision by the campaign
owner; the joint critic's disposition is the same, recording the fence as a
scope statement and a lane record rather than a claim. What did travel from
this lane into the claims DAG is only the conditional selection theorem on the
two-dimensional register, which carries this fence in its own scope.

\provenance{no claims-DAG row (negative lane record and unpromoted companion material); the fence is carried inside A-INDEX-PEPS's not-claimed list}%
{\texttt{theory/soft-2d-hunt.md} §3 (the no-go and its probe-family fence) and §4 (the exact stabiliser patch); \texttt{theory/cat-hunt.md} §§2--3, §5.3 (the operator/categorical split) and §9 (honest status)}%
{\texttt{theory/checks/cat\_hunt\_check.py} CATH-C0--C4 with failing mode \texttt{-{}-red missing-ising-channel}; \texttt{theory/checks/soft\_2d\_hunt\_check.py} S2DH gates}%
{--- (the joint critic \texttt{theory/verdicts/w21-joint-critic.md} W6-O2 fenced the statement's strength and directed that it not be promoted; it adjudicated the neighbouring positive rows, not this negative one)}

\subsection{What the negatives have in common}
\label{sub:negative-synthesis}

Read together, the ledger and the three obstructions show one pattern, and it
is worth naming because it predicts where the next dead end will be. Every
failed universality in this campaign was a statement trying to obtain a
\emph{number} from symmetry alone. The withdrawn orbit label wanted a complete
invariant from a group quotient, and got one only after the diagonal action was
taken seriously and the answer turned out to be a double coset. The withdrawn
soft factor wanted a hard-independent coefficient from a kinematic difference,
and lost it as soon as one asked what sits inside the current matrix element.
The unrestricted universality claim wanted one coefficient for every local
source, and a four-site operator moved that coefficient while leaving the hard
process fixed. The literal memory conjecture wanted a displacement from a
phase, when the displacement is a charge. The amputation conjecture wants a
value from a normalisation, and normalisation supplies exactly one square root
of it. The localisation hunt wanted outcome probabilities from a volume, and
volumes consume multiplicities rather than producing them. The categorical hunt
wanted a Taylor coefficient from fusion, and fusion is blind to any scalar
dressing that is trivial at zero momentum.

The surviving architecture is the complement of that list, and it is exactly
what the positive chapters prove. Symmetry — group-theoretic, categorical, or
cohomological — delivers structure: which sectors exist, which transitions are
allowed, which observables are quantised, and in what coset their spectrum
lies. Dynamics delivers values: the two-body contact equation gives the soft
slope $1/S$, the current and channel velocity give the residue $2iv_h$, and the
scattering problem gives the transmission weights. Every proved statement in
this labbook respects that division of labour, and every refuted one violated
it. The practical rule extracted from the negatives is therefore simple: when a
proposed theorem promises a number, one should be able to point at the
dynamical input that carries it — and if the only inputs are a group, a
category, a class definition or a volume, the number is not there.
